\section{When an end-to-end graph workflow still lacks calibrated effects}
\label{rel64:graph-confidence}

The preceding example leaves a practical question. Can an assistant quantify
its uncertainty about each arrow, draw many graphs and pass those graphs to a
flexible generative model? Such a workflow would be useful if its final
intervals retained the mechanisms relevant to the bank's decision. Each step,
however, changes the object whose uncertainty is being described. A probability
assigned to a sentence, a distribution over graphs and a sampling interval for
an intervention effect require different justifications.

Molero P{\'e}rez's February 2026 thesis develops a concrete instance of this
workflow. Pairwise language judgments, optionally augmented by retrieval,
produce proposed causal relations. The study compares self-reported
confidence, selected-answer token probabilities and a HallBayes refinement
based on weakened prompts. Its chosen procedure converts the refinement's
risk score into an edge probability, averages three runs, samples 100 graphs
and fits causal neural spline flows. This is a substantive implementation
proposal with reported biomedical comparisons. Its availability permits a
more precise assessment than either accepting the integration as calibrated
or dismissing it merely because language models propose the graphs. The
following assessment concerns the printed method and results; it does not
reproduce the upstream software or its model calls.\citep{rel64molero}

\subsection{A success ceiling gives an error floor}

Write the Bernoulli relative entropy, in nats, as
\[
 d(p,q)=p\log\frac pq+(1-p)\log\frac{1-p}{1-q},
 \qquad 0<q<1.
\]
The thesis's equations (14)--(16) average clipped positive log-probability
lifts for the selected answer under original and masked prompts. It then uses
this quantity as an information budget in a Bernoulli inequality and defines
the largest compatible success probability. For any stipulated nonnegative
budget $\Delta$, the logical implication of that inequality is
\begin{equation}
 \begin{split}
 d(p,q)\leq\Delta,\qquad
 p_{\max}(\Delta,q)&=\sup\{u\in[0,1]:d(u,q)\leq\Delta\},\\
 p\leq p_{\max}(\Delta,q)&\quad\Longrightarrow\quad
 1-p\geq1-p_{\max}(\Delta,q).
 \end{split}
 \label{rel64:error-floor}
\end{equation}
The complement is therefore a lower bound on error under this premise. The
thesis calls it an upper bound and subsequently uses it to set graph-sampling
probabilities. That interpretation does not follow from the displayed
inequality. With $q=.1$ and $\Delta=.5$, the upper endpoint is approximately
$.49505$, giving complement $.50495$. Yet $p=.1$ satisfies the inequality
and has error $.9$. This exact counterexample concerns the direction of the
bound; it does not assert that every estimated graph in the thesis is wrong.
\citep{rel64molero}

There is a further distinction between the score and the premise. Evaluating
a clipped log-lift at one selected answer does not generally evaluate an
expectation over the full answer distribution. Nor does masking a prompt
establish the probability that its causal assertion is true in the intended
bank population. A separate calibration argument is needed before the score
can stand for the information budget or the causal correctness probability
in~\eqref{rel64:error-floor}. Averaging repeated calls reduces some
Monte Carlo variation conditional on the response mechanism; it cannot by
itself supply either missing interpretation.

\subsection{The information inequality remains useful in its stated direction}

The thesis cites the 2025 version of Chlon and colleagues' information-budget
paper. Its Theorem 5 follows from convexity and reduction to an event. Let
$P$ be a reference distribution, $S_k$ predictive distributions, and
$\bar S=\mathbb E_k S_k$ their mixture. For an event $A$, put
$p=P(A)$ and $q=\bar S(A)$. Under absolute continuity and finite expected
relative entropy,
\begin{equation}
 \mathbb E_k\operatorname{KL}(P\Vert S_k)
 \geq\operatorname{KL}(P\Vert\bar S)
 \geq d(p,q).
 \label{rel64:necessary-information}
\end{equation}
The first step is Jensen's inequality in the second distribution; the second
is the information lost when the outcome is reduced to membership of $A$.
Equality for the entire chain requires equality in both steps. Sufficiency
of the event for distinguishing $P$ from the mixture addresses the second
step; it does not eliminate a positive mixture Jensen gap.
\citep{rel64chlonone,rel64chlonthree}

A necessary information requirement is not a success guarantee. Consider the
project's four-outcome construction
\[
 S=(.05,.05,.001,.899),\qquad
 P=(.05,.05,.899,.001),\qquad A=\{1,2\}.
\]
Use the same $S$ for every prompt. Both event probabilities are $.1$, but
\[
 \operatorname{KL}(P\Vert S)=.898\log899\simeq6.10755
 >d(.95,.1)\simeq1.99421.
\]
The extra divergence lies within the failure event. Exceeding the information
requirement associated with a target success probability of $.95$ has not
raised actual reference-event probability above $.1$. This is a failure of a
sufficiency interpretation, not a counterexample to
inequality~\eqref{rel64:necessary-information}.

The June 2026 version changes the title, author list and deployment
description. It explicitly requires a specified Bernoulli predicate and a
reference supplied by gold labels, a deterministic verification primitive or
an audited calibrated verifier. Unrestricted generation without such a
reference is outside its stated regime. Its reported held-out audit remains
one particular evaluation, not an intervention-coverage result. These
qualifications improve the interpretation of the proposal, but they do not
turn the necessary inequality into a sufficient guarantee. Even with a
binary reference, $d(.05,.5)=d(.95,.5)$: without a justified branch
restriction, inversion does not distinguish low from high event probability.
With constant predictive probability $.5$ and target $.95$, the printed
budget ratio is one for either reference probability, before clipping matters.
The event, reference and direction of the bound must therefore remain visible
when adapting the planner. A later version's verifier assumptions also cannot
be silently attributed to the thesis's earlier masking procedure.
\citep{rel64chlonthree}

\subsection{Independent edge draws need a coherent graph law}

For an accepted edge the thesis uses one minus the risk score; for a rejected
edge it uses the risk score. It then samples edges independently and applies
topological sorting. Sorting represents an already acyclic graph in an
appropriate order; it cannot make every independently drawn graph acyclic.
The treatment of a failed sort is part of the probability law. The following
example diagnoses the general construction and does not claim that the
thesis's particular draws contained a cycle.\citep{rel64molero}

Suppose the only possible arrows are $A\to B$, $B\to C$ and $C\to A$,
each assigned probability $4/5$. Under independent draws,
\[
 \Pr(\text{cycle})=(4/5)^3=64/125,\qquad
 \Pr(A\to B\mid\text{acyclic})=\frac{36}{61}\simeq.59016.
\]
The conditional numerator is $(4/5)[1-(4/5)^2]$ and the denominator
is $1-(4/5)^3$. Rejecting cyclic draws changes the edge probabilities;
deleting an edge instead creates another distribution determined by the
deletion rule. Indeed, no distribution supported on DAGs can retain all
three proposed marginals: a DAG contains at most two of these arrows, whereas
their proposed expected count is $12/5$. This strengthens the majority-cycle
example in Section~\ref{exec20:expert-votes}: structural validity and
marginal probability interpretation cannot be checked independently.

A usable alternative defines a joint law on an explicitly admissible model
class, records its prior and likelihood interpretation, and tests sensitivity
to that class. When those probabilities cannot be justified, retain compatible
graphs and propagate their effect sets instead. Neither route assumes that a
linguistic probability score is already a posterior over banking mechanisms.

\subsection{An interval comparison needs an intervention reference}

The thesis trains four groups of 100 flow models: an expert reference graph,
the sampled language-derived graphs, the expert graph's ordering without its
adjacency restrictions, and random orderings without those restrictions.
Within each group, the 2.5th and 97.5th percentiles of the fitted effect
estimates define an interval. The reported interval-overlap means are .723,
.712 and .301 for the last three groups compared with the fitted expert-graph
group. The first two differ by .011. That contrast describes these fitted
model intervals; the table does not provide a paired uncertainty interval for
the difference. It is not evidence that language assistance generally
outperforms knowledge of the causal ordering.\citep{rel64molero}

Agreement with a fitted expert reference is useful when the task is to
reconstruct that reference. It is not observed intervention truth. If two
methods return the identical interval $[.5,.7]$ while the true effect is zero,
their interval overlap is one and both miss the effect. More generally,
percentiles across graph choices and training runs describe variation under
that procedure. A 95\% repeated-sampling coverage claim needs a defined
sampling experiment, valid treatment of graph uncertainty and an identified
target. A calibrated Bayesian interpretation instead needs a justified joint
posterior and corresponding model assumptions. Repetition alone supplies
neither interpretation.

The numerical comparator also needs matched information and a suitable test.
The thesis specifies chi-square conditional-independence tests for its PC
baseline while prioritizing continuous clinical variables. Its language
workflow additionally draws on outside knowledge. A fair follow-up would
resolve the continuous-variable testing and coding choices, compare methods
given the same admissible knowledge, and keep a numerical-only arm. The
thesis itself acknowledges calibration problems, lost valid edges and the
absence of hidden-confounder handling. Those limitations belong alongside
the favorable reported scores.\citep{rel64molero}

For the bank, the investment conclusion follows from the retained example in
Section~\ref{exec20:expert-identification}. Its observationally identical
mechanisms imply margins of USD~20 and USD~$-40$ under the same stated
valuation and cost. A narrow distribution of repeated fits cannot resolve
that disagreement. Before releasing a financial recommendation, the workflow
must state the decision-time information that excludes a mechanism, or carry
the ambiguity into the recommendation. Graph proposals, retrieval and flexible
conditional models remain candidate components. Their integration deserves
evaluation against an expert-reviewed graph, a numerical discovery route and
an ambiguity-preserving analysis, with withheld mechanisms, interval coverage,
decision loss and complete review cost assessed separately.
